\section{Discussion}
\label{sec:discussion}

\paragraph{Progressive does not mean flat.}
Joggle programs retain functions, blocks, operations, values, and def--use
relations. These structures express program semantics and support conventional
local reasoning. Progression changes how stages relate: compiler functions
refine a verified graph without requiring a global ladder of public IR classes.
High-level operations and lower-level helpers may therefore coexist when a
stage needs both.

\paragraph{A mod is a capability boundary.}
Directories split files, and hierarchical IR splits programs; neither
necessarily identifies the owner of a compiler feature. A mod binds a public
typed surface to dependencies, source fragments, native bindings, and lifecycle
operations. Small implementations may remain in one source file. Separate mods
become useful when capabilities require distinct owners, dependency closures,
or release cycles.

\paragraph{Uniform calls preserve effects.}
Uniform syntax coexists with explicit effect contracts. Query execution is
read-only, runs publish a verified graph, and artifact calls return owned
values. Native implementations use the same signatures and ownership boundary.
Declarations, resolution, calls, and composition are uniform; effect checks
remain explicit.

\paragraph{Precision has a cost.}
A compiler function gains reuse by reading the narrowest state needed for its
result. Whole-graph traversal correctly creates a broad dependency; range and
single-entity access create narrow ones. Narrow records add capture and
validation work, so they help only when avoided execution is more expensive.
Miss reasons, observation counts, and executed stage counts expose this
trade-off directly.

\paragraph{Publication bounds reuse.}
Dependency records are valid only for the graph that produced them. Joggle
therefore commits mutations, revisions, and new records after final
verification, or restores the previous state and discards tentative records.
State outside the graph must enter through typed arguments or an environment
revision. A later run can then reuse only observations tied to a published
graph and environment.

\paragraph{Plans remain interpreted.}
Predecoded plans cache checked interpreter work: dense value slots,
control-flow targets, operand indices, and call-site information. They remove
repeated setup while preserving interpreted execution. Native acceleration
remains a typed mod binding, and native plan compilation is an independent
design point.
